
\begin{figure}
\begin{minipage}{1\columnwidth}
    \centering
\begin{tcolorbox}[title=Relation Selection Prompt]
        \small
        \vspace{10pt}
        \textbf{Task: Relation Selection} \\
Your task is to select the most appropriate relations based on their relevance to a given question.  \\



    Follow these steps: \\
    1. Read the provided <question>question</question> carefully. \\
    2. Analyze each relation in the <relation> list and determine its relevance to the question and relation. \\
    3. Respond by selecting the most relevant relations within <select>relations</select> tags. Be both frugal on your selection and consider completeness. \\
    4. The selected relations should be provided line-by-line. \\
    

    Example format:
    <question> \\
    Name the president of the country whose main spoken language was English in 1980? \\
    </question> \\

    <entity> \\
    English \\
    </entity> \\
 
    <relations> \\
    language.human\_language.main\_country \\
    language.human\_language.language\_family \\
    language.human\_language.iso\_639\_3\_code \\
    base.rosetta.languoid.parent \\
    language.human\_language.countries\_spoken\_in \\
    </relations> \\

    <selected> \\
    language.human\_language.main\_country \\
    language.human\_language.countries\_spoken\_in \\
    base.rosetta.languoid.parent \\
    </selected> \\

    Explanation: [short explanation (deprecated due to figure length)] \\

    Important Instructions: Always return at least one relation.
    Now it is your turn. 

    \vspace{5pt}
    <\textbf{question}> \\
    \{question\} \\
    </\textbf{question}> \\

    <\textbf{entity}> \\
    \{entity\}  \\
    </\textbf{entity}> \\
 
    <\textbf{relations}> \\
    \{relations\} \\
    </\textbf{relations}> \\

    Remember to parse your response in <selected></selected> tags:
    \end{tcolorbox}

\caption{The relation selection prompt template used in agentic traversal. }
\label{fig:agent-relation-prompt}
\end{minipage}
\end{figure}
